\section{Native-oriented Artin--Schreier classes}\label{sec:as}

The field result permits a concrete, sign-fixed resolvent.
For \(n=p^e\) and \(\alpha_e\in S_e\), set
\begin{equation}\label{eq:resolvent}
 w_e=\frac{\alpha_e^{n-1}}{M_e'(\alpha_e)},\qquad
 y_e=-\sum_{i=0}^{n-1}\overline i\,\sigma_e^i(w_e),\qquad
 a_e=y_e^p-y_e ,
\end{equation}
where \(\overline i\in\Fp\) is the residue of \(i\).

\begin{lemma}[Positive native translation]\label{lem:resolvent}
The elements in \eqref{eq:resolvent} satisfy
\[
 \Tr_{L_e/K}(w_e)=1,\qquad
 \sigma_e(y_e)-y_e=1,\qquad a_e\in K,\qquad K(y_e)=F_e.
\]
The class \([a_e]\in K/\wp(K)\) is independent of the starting
root and of the trace-one choice, with this native orientation.
\end{lemma}

\begin{proof}
Let \(x_0,\ldots,x_{n-1}\) be the distinct roots of \(M_e\).
The coefficient of \(T^{n-1}\) in Lagrange interpolation
\[
 T^{n-1}=\sum_{i=0}^{n-1}
 x_i^{n-1}\frac{M_e(T)}{(T-x_i)M_e'(x_i)}
\]
gives \(\sum_i x_i^{n-1}/M_e'(x_i)=1\).
This is the trace identity. Reindexing the weighted sum
in \eqref{eq:resolvent}, including its wrapped coefficient
\(-(n-1)=1\) in characteristic \(p\), gives
\(\sigma_e y_e-y_e=\sum_i\sigma_e^i w_e=1\).
Thus \(\sigma_e^a y_e=y_e+\overline a\), \(a_e\) is invariant,
and the stabilizer of \(y_e\) is precisely
\(\langle\sigma_e^p\rangle\). This proves \(K(y_e)=F_e\).

Replacing \(w_e\) by any trace-one element produces \(y'_e\)
with the same translation identity. Their difference is fixed
by \(\sigma_e\), hence lies in \(K\), and
\(a'_e-a_e=\wp(y'_e-y_e)\). Replacing the starting root
by \(\sigma_e^a\alpha_e\) replaces \(y_e\) by
\(\sigma_e^a y_e=y_e+\overline a\), also preserving the class.
No division by \(p^e\) has been used.
\end{proof}

This is the additive trace-resolvent construction of classical
Artin--Schreier theory, written here to specify the native sign.
For completeness, every class in \(K/\wp(K)\) has a unique
reduced polar representative
\(\sum_{m>0,\ p\nmid m}c_ms^{-m}\) with finite support.
The regular part is in \(\wp(k[[s]])\): the constant equation
is solvable in \(k\), and the remaining coefficients are solved
successively using the unit derivative \(-1\).
Subtracting \(\wp(c^{1/p}s^{-m})\) eliminates a pole
\(cs^{-pm}\); starting from the largest pole terminates.
A nonzero reduced polar polynomial cannot equal \(\wp(b)\):
if \(v(b)<0\) its largest pole would be divisible by \(p\),
and if \(v(b)\ge0\) it would have no pole. This also proves
uniqueness.

\begin{proposition}[Class stabilization]\label{prop:as-stable}
For every \(e\ge2\), \([a_e]=[a_2]\) and \(F_e=F_2\)
inside the fixed separable closure.
\end{proposition}

\begin{proof}
Let \(\rho_e:G_K\to\Z/p^e\Z\) be characterized by
\[
 g\alpha_e=P^{\circ\rho_e(g)}(\alpha_e).
\]
Commutation with \(P\) shows that the index is independent of
the starting root; composition adds indices, so it is a
continuous character, factoring through \(L_e/K\).
For \(e\ge3\), label the top clusters at each level by
\(i\in\Z/p\Z\), using native iteration.
By Lemma~\ref{lem:matching}, the matching is a map commuting
with \(i\mapsto i+1\), hence has the form \(i\mapsto i+c\).
Galois equivariance then gives, for every \(g\in G_K\),
\[
 i+\overline{\rho_e(g)}+c
 =i+c+\overline{\rho_2(g)} .
\]
Thus \(\rho_e\bmod p=\rho_2\bmod p\) on the whole group.
The same assertion is tautological at \(e=2\).
The translation identity in Lemma~\ref{lem:resolvent} gives
\[
 g(y_e)-y_e=\overline{\rho_e(g)}
 =\overline{\rho_2(g)}=g(y_2)-y_2.
\]
Consequently \(b_e=y_e-y_2\) is fixed by \(G_K\) and lies
in \(K\). It follows that \(a_e-a_2=\wp(b_e)\), and that
\(K(y_e)=K(y_2)\).
The phase \(c\) permits only a translation of cluster labels,
not a nontrivial scalar or sign change of the native generator.
\end{proof}

\begin{proposition}[Separation of the prime-period field]\label{prop:separation}
For every \(e\ge2\), \(L_1\cap L_e=K\).
\end{proposition}

\begin{proof}
First suppose, for contradiction, \(L_1\subset L_2\).
Choose \(\alpha\in S_2,\beta\in S_1\), and put
\(\eta=\alpha-\beta\in L_2\).
By Proposition~\ref{prop:second},
\[
 v_{L_2}(\eta)=2(p-1)^2,\qquad
 v_{L_2}(\sigma_2^p\eta-\eta)=2p^2(p-1).
\]
For the second equation, \(\sigma_2^p\) fixes the unique
degree-\(p\) subfield, hypothetically \(L_1\), and the
\(p\)-step displacement of \(\alpha\) is \eqref{eq:pstep}.
The first value is prime to \(p\), so
Lemma~\ref{lem:prime-value} gives the second lower break,
under this supposition, as
\[
 b_2=2p^2(p-1)-2(p-1)^2
 =2(p-1)(p^2-p+1).
\]
The degree-\(p\) quotient is hypothetically \(L_1/K\), giving
first break \(b_1=p-1\). But then
\[
 b_2-b_1=(p-1)(2p^2-2p+1)\equiv-1\pmod p,
\]
contrary to \eqref{eq:herbrand-two}.
Thus \(L_1\not\subset L_2\), and its prime degree gives
\(L_1\cap L_2=K\).

If \(L_1\subset L_e\) at any later level, cyclicity would
identify it with \(F_e=F_2\subset L_2\), contradicting what
was just proved. Again prime degree gives the stated
intersection. This completes Theorem~\ref{thm:main}.
\end{proof}

The second lower-break expression in this contradiction has
not yet been proved to be an actual invariant: its paired
first break arose from a false containment assumption.
Section~\ref{sec:ramification} derives the actual two breaks
independently, with first break \(2(p-1)\).
